\documentclass{article}

% Language setting
\usepackage[english]{babel}
\usepackage{siunitx}

% Set page size and margins
\usepackage[letterpaper,top=2cm,bottom=2cm,left=3cm,right=3cm,marginparwidth=1.75cm]{geometry}

% Useful packages
\usepackage{amsmath}
\usepackage{graphicx}
\usepackage[colorlinks=true, allcolors=blue]{hyperref}

\title{TDT4173 Course Project}
\author{Nymoen, Birkeli}

\begin{document}
\maketitle

\noindent
\textbf{Group:} 61 \quad \\
\textbf{Kaggle team:} [61] Nymoen Birkeli \\
\textbf{Members:} Fredrik Nymoen, student ID [FILL IN] \; Jonas Birkeli, 107006

\begin{abstract}
We predict unit commitment for the Tokke--Vinje hydropower system: for
every daily scheduling case (a 168-hour week), a continuous ON-probability
for each of 14 generators and 168 hours, graded by micro-averaged ROC-AUC.
The labels are the output of a scheduling optimizer, so the task is to
imitate its strategy from its own information state. We built two model
families on a shared physical feature inventory: 14 per-generator XGBoost
classifiers (validation micro AUC 0.9857 on 2022, public leaderboard
0.99066) and one shared temporal convolutional network with local plant
branches (validation micro AUC 0.9809/0.9893 on purged 2020/2021 folds,
public leaderboard 0.99261). The report documents the full search path,
including the negative results: closed feature-search rounds, a closed
hyperparameter screen, and rejected model families.
\end{abstract}

\section{Introduction}

Unit commitment in a hydropower cascade is a week-level regime decision
with hourly modulation, driven by the comparison between the electricity
price and the value of the water behind the dam. The labels come from a
scheduling optimizer re-solved daily, so the target is a \emph{strategy},
not a physical observation, and the same physical hour carries different
labels in overlapping cases.

Exploratory analysis showed that water state, not price level, sets the
operating regime. We encoded that physics explicitly, price minus water
value, water-system allocation, terminal water-value conversion and
start/stop-cost scenarios, and fed the resulting features to two model
families: per-generator gradient-boosted trees and a shared temporal
convolutional network with local plant branches. We fixed a leak-free
chronological validation protocol before comparing anything. The report
documents what did work and what did not.

\section{Task and Data}

\subsection{Data}
\textbf{Inputs.} Five input files: hourly day-ahead price (2015--2025),
daily reservoir volumes and daily synthetic water values for 17 reservoirs
(2015--2024), hourly per-reservoir inflow back to 1958, and the unit
commitment decisions, 2922 daily cases (2015--2022), each with 14
generators $\times$ 168 hours of binary labels. Extended data provides
minimum-flow and minimum-volume constraints and a topology file mapping
plants to their supplying reservoirs with level--volume curves.

\textbf{Prediction target.} For each of the 725 hidden test cases
(2023-01-01 to 2024-12-25), a probability in $[0,1]$ for every
(case, hour, generator) point, 1.7M predictions per submission.

\subsection{Metric}
\label{sec:metric}
Micro ROC-AUC pooled over all (case, hour, generator) points. This has a
practical consequence we measured directly: per-generator rank
normalization destroys the cross-generator score levels and collapses the
metric (a rank ensemble scored 0.9556 vs.\ 0.9809 for a mean-probability
ensemble), so all ensembling is done on raw probabilities.

\subsection{Validation Protocol}
One rule was used everywhere: \textbf{train on earlier years, score on
later years, and never let the two overlap}. Each case covers a whole
week, and the same physical hour can have different labels in different
overlapping cases (Section~\ref{sec:intuitive}), so every training week
must end before the scoring period begins. Any statistics the model learns
from the data (scaling, normalization) are computed on training cases
only.

Two setups were used under this rule:
\begin{itemize}
    \item \textbf{Tree models:} train on 2015--2021, score on 2022. The
    model decides when to stop training using only the last 90 days of
    the training period, never the scoring year. An earlier version
    stopped training by looking at the scoring year itself, which leaked
    information and made results swing by $\pm 0.02$; fixing this was
    worth $+0.012$ AUC (0.9686 $\to$ 0.9802). A new idea was only kept
    if it improved the score by more than the run-to-run noise
    ($+0.0015$; the measured noise floor is $\pm 0.0014$).
    \item \textbf{Sequence model:} train on 2015--2019 and score on 2020,
    then train on 2015--2020 and score on 2021, with training weeks
    trimmed so nothing overlaps the scoring year. The number of training
    epochs was fixed in advance. The final submission retrains on all
    2922 cases.
\end{itemize}
As a stability check we also ran leave-one-year-out cross-validation:
pooled out-of-fold micro AUC 0.9758 over all 2922 cases (the 2020
COVID-outlier fold is the weakest at 0.9018).

\section{Exploratory Data Analysis}
\subsection{Search Domain Knowledge}
\label{sec:domain}

A hydropower generator is committed when producing electricity now is
worth more than storing the water. The scheduling model therefore
compares, for every hour, the \textbf{day-ahead electricity price}
against the \textbf{water value} of each reservoir. The opportunity cost
of releasing one more unit of water, which depends on how full the
reservoir is and on expected future inflow and prices. This makes
\emph{price minus water value} the natural decision variable, and it is
provided per reservoir. Three further physical effects matter:
\begin{itemize}
    \item \textbf{Head}: the power produced per
    \SI{}{\cubic\meter\per\second} grows with the height difference
    between the reservoir surface and the turbine, so the reservoir
    \emph{level} (not just the volume) affects how profitable a generator
    is at a given moment.
    \item \textbf{Cascade coupling}: the system is a chain of reservoirs
    and plants. Water released upstream reappears in a downstream
    reservoir, so the commitment of one plant depends on the state of the
    reservoirs above it, not only its own.
    \item \textbf{Constraints}: environmental minimum-flow and
    minimum-volume rules can force generation (or spill) regardless of
    price, decoupling commitment from the price signal for some units.
\end{itemize}
Inflow drives everything on longer time scales: in a wet year reservoirs
fill, water values fall and units run even at low prices. This domain
picture directly motivated our key feature family, price minus water
value, and the per-plant aggregates over upstream reservoirs described in
Section~\ref{sec:features}.

\subsection{Check If The Data Is Intuitive}
\label{sec:intuitive}

Mostly yes, with two exceptions.
\begin{itemize}
    \item \textbf{Intuitive}: commitment is strongly sticky. Hour-to-hour
    transition probabilities are P(on$\mid$on) 0.96--0.997 and
    P(off$\mid$off) 0.97--0.997 per generator, with hourly flip rates of
    only 0.3--2.2\%. Between 30\% and 80\% of case-weeks are purely
    all-on or all-off (Haukeli 81\%): the label is mostly a week-level
    state decision with hourly modulation. Commitment also follows the
    expected daily and annual cycles (ON-rate by hour of day and by day
    of year, Figure~\ref{fig:cycles}).
    \item \textbf{Counter-intuitive at first}: ON-rate is non-monotone in
    price. Most units are most ON in the lowest price domain. Exactly
    what the domain picture predicts (spring flood: full reservoirs, low
    prices, cheap water), but the opposite of a naive price-threshold
    model. Water state dominates price level.
    \item \textbf{Counter-intuitive at first}: the same physical hour
    carries different labels in the (up to 7) overlapping cases that cover
    it. The per-generator mismatch rate ranges from 4.9\% (Byrte) to
    18.7\% (Tokke\_G2). This is consistent with how the data is generated
    (Section~\ref{sec:generated}): each case is solved independently, so
    the commitment of an hour depends on the case horizon. It also
    explains why pooling predictions across overlapping cases is wrong,
    and why the in-sample AUC ceiling is 0.9985 and not 1.0.
\end{itemize}

\subsection{Understand How the Data Was Generated}
\label{sec:generated}

The labels are not measurements of what the plants did, but rather the
output of a short-term stochastic scheduling optimizer, re-solved
independently for every daily case. Each solve takes the case-start
reservoir levels, the hourly day-ahead price, hourly inflow and daily
synthetic water values as input, and commits the 14 generators for the 168
hours of that week. The provided inputs are therefore exactly the
optimizer's information state, and the learning task is to imitate the
optimizer's strategy. Two consequences we verified:
\begin{itemize}
    \item The water values are synthetic (model-generated, not market
    prices), and the inflow history extends back to 1958, which allows
    climatological baselines.
    \item The optimizer's own difficulty varies. The per-case solve time
    has a median of 52\,s but a mean of 147\,s and a maximum of 4526\,s, a
    small share of cases are genuinely hard, and those are the ones where
    the strategy is hardest to imitate.
\end{itemize}

\subsection{Explore Individual Features}
\label{sec:individual}
\paragraph{Price.} Hourly day-ahead price, 2015--2025, mean 59.5 EUR/MWh,
range $-61.8$ to $898.3$. Negative prices occur in 0.44\% of hours. But,
only 10 negative hours fall before 2023, so negative prices are
essentially a test-period phenomenon and a real train/test distribution
shift.

\paragraph{Reservoir state and water values.} Daily volumes for 17
reservoirs (2015--2024) and a synthetic water value per reservoir per day.
The derived column \emph{price minus water value} is the single most
predictive feature family: in a per-generator univariate screen
(direction-corrected AUC), the best single column,
\texttt{price\_minus\_val\_\_Vinjevatn}, reaches a mean AUC of 0.826
across the 14 generators, and the entire top-20 is dominated by
\texttt{price\_minus\_val\_\_*} and terminal-water-value columns.

\paragraph{Inflow.} Hourly per-reservoir inflow back to 1958. In the
labeled period it is genuinely hourly; in the early decades the file
repeats daily values on the hour (Section~\ref{sec:cleanup}).

\paragraph{Calendar.} ON-rate panels by hour of day, month and day of
year show the expected daily and annual cycles, which motivates the
sin/cos encodings of hour, day-of-week and day-of-year used in the feature
set.

\paragraph{Labels.} Strictly binary $\{0,1\}$, no missing values.
Overall ON-rates range from $0.26$ (Vinje\_G3) to $0.75$ (Lio\_G1) and
drift strongly by year, e.g.\ Haukeli falls from $0.87$ (2015) to $0.14$
(2022) and Songa from $0.70$ (2020) to $0.12$ (2022). So the test years
sit at the end of a decade-long drift, not in the middle of it.

\subsection{Explore Pairs and Groups of Features}
\label{sec:pairs}

\textbf{Price vs.\ water state (the decisive pair).} Tabulating ON-rate
per price decile per generator shows the non-monotone price response
described above: water state, not price level, sets the operating regime.
This is confirmed by a feature-group ablation (default XGBoost, train
2015--2021, val 2022, micro AUC):

\begin{table}[h]
    \centering
    \begin{tabular}{lrr}
        \hline
        Configuration & Micro AUC & $\Delta$ \\
        \hline
        All features & 0.9764 & --- \\
        $-$ reservoir / water-value group & 0.9355 & $-0.0410$ \\
        $-$ price group & 0.9729 & $-0.0035$ \\
        $-$ calendar group & 0.9757 & $-0.0007$ \\
        $-$ inflow group & 0.9767 & $+0.0003$ \\
        $-$ constraints group & 0.9776 & $+0.0012$ \\
        $-$ local aggregates & 0.9781 & $+0.0017$ \\
        Top-10 features only & 0.9548 & $-0.0216$ \\
        \hline
    \end{tabular}
    \caption{Feature-group ablation. The reservoir/water-value group
    carries nearly all the signal; no single small set of features
    suffices.}
    \label{tab:ablation}
\end{table}

\textbf{Signal is spread thin.} A model restricted to the top-10
univariate features drops to 0.9548. There is no small magic subset,
which argues for a model class (gradient boosting) that can combine many
weak features.

\textbf{Hard units respond to context, but in opposite directions.}
Features describing upstream water state helped Haukeli ($+0.006$) and
Vesle Kjela ($+0.007$) but hurt Lio ($-0.003$), netting to zero at the
pooled level. This recurring see-saw pattern across five feature rounds
is what closed the feature search: the two hard units fail for different
reasons (Section~\ref{sec:interpretation}).

\subsection{Cleanup Features}
\label{sec:cleanup}

We checked every input file systematically:
\begin{itemize}
    \item \textbf{No missing values and no duplicate timestamps} in any
    of the six input files (price, volume, water value, inflow, minimum
    flow, minimum volume). No imputation was needed.
    \item \textbf{Labels are strictly binary} $\{0,1\}$ with no
    corruption; they are cast to float only for training.
    \item \textbf{Resolution alignment}: price and inflow are hourly,
    volumes and water values are daily. Daily values are aligned to the
    case start, and per-hour features are looked up by timestamp. The
    inflow file is hourly only in the modern period; before that, daily
    values are repeated on each hour (154{,}879 consecutive identical
    hourly rows overall vs.\ only 50 in 2015+), which matters only if
    pre-2015 inflow is used (we use it solely as a climatology baseline,
    where the repetition is harmless).
    \item \textbf{Price resolution}: 15-minute prices exist only from
    October 2025, after the test window, so hourly is the finest usable
    resolution.
    \item \textbf{Dropped columns}: the optimizer solve-time column
    exists only for training cases and cannot be used as a feature; the
    topology file is used only to derive static features
    (plant-to-reservoir mapping, level curves), not as a data table.
    \item No outlier removal was applied: extreme prices (up to 898
    EUR/MWh) and negative prices are real operating conditions the model
    must handle.
\end{itemize}

\section{Feature Engineering}
\label{sec:features}

The feature inventory is shared by both model families. For each hour of
a case-week:
\begin{itemize}
    \item \textbf{Calendar} — hour of day, day of week, month, day of
    year (as sin/cos curves), and how far into the week we are.
    \item \textbf{Price} — the hourly price, weekly average/min/max,
    rank of this hour within the week, averages of the first/last hours.
    \item \textbf{Reservoir state} — how full each reservoir is at the
    start of the week, and the water value at the end of the week.
    \item \textbf{Price minus water value} — the key family. If
    electricity is expensive but water is cheap to use, the generator
    should be ON. This one family carries almost all the signal.
    \item \textbf{Inflow} — how much water flows into each reservoir,
    cumulative and remaining over the week.
    \item \textbf{Constraints} — minimum flow / minimum volume rules.
    \item \textbf{Local aggregates} — the same quantities, averaged over
    just the reservoirs feeding this specific plant.
    \item \textbf{LP water-system allocation} — a simplified linear
    program routes water through the named topology and returns
    projected plant flows, reservoir storage and marginal water costs.
    Used by both families.
    \item \textbf{Terminal water-value conversion} — local terminal
    EUR/MWh values converted through best-operating-point energy at
    initial head and downstream receiving-reservoir values.
    \item \textbf{Start/stop-cost scenarios} — 54 binary channels
    estimating the whole-week payoff of staying ON through a bad hour
    (9 plants $\times$ 3 cost scales $\times$ 2 initial states).
    \item \textbf{Asinh encoding} — extreme prices and water values are
    squashed with asinh instead of clipped, keeping their sign and
    magnitude information.
\end{itemize}

The feature search was run one factor at a time against a fixed anchor
under the adoption bar. Table~\ref{tab:arms} shows the final round: each
arm adds one feature group to the anchor configuration.

\begin{table}[h]
    \centering
    \begin{tabular}{llrr}
        \hline
        Arm & Added features & Val 2022 & $\Delta$ \\
        \hline
        anchor & v2 + 12h LP allocation & 0.9842 & --- \\
        lp2 & LP duals + min-volume margin & 0.9840 & $-0.0002$ \\
        margin\_ctx & profitable-run length and margins & 0.9850 & $+0.0008$ \\
        duals & LP marginal water cost & 0.9854 & $+0.0012$ \\
        hydraulics & tunnel submergence, spillway capacity & 0.9855 & $+0.0013$ \\
        \textbf{minvolmargin} & projected margin to min-volume floor & \textbf{0.9857} & $+0.0015$ \\
        \hline
    \end{tabular}
    \caption{Feature arms, each one factor apart from the anchor (same
    protocol, same hyperparameters). The winning arm (minvolmargin: the
    projected reservoir margin to the minimum-volume floor along the LP
    storage trajectory, where the baseline only had the margin at case
    start) reaches the adoption bar and was retrained on all 2922 cases
    for the final submission.}
    \label{tab:arms}
\end{table}

What never helped: more price history, inflow anomalies, or anything
encoding the same physics a second way. The feature search was closed
after five separate rounds of null results.

\section{Models}

\subsection{Gradient-Boosted Trees}
\label{sec:xgboost}

\textbf{Setup.} One XGBoost classifier per generator: 14 models, each
trained on its own generator's labels. Each model sees
$2{,}557 \times 168 = 429{,}576$ training rows (2015--2021 case starts)
and is scored on $365 \times 168 = 61{,}320$ rows (2022).

\textbf{Features.} 193 columns per row, dominated by the \emph{price
minus water value} family and per-plant aggregates
(Section~\ref{sec:features}). The feature-group ablation
(Table~\ref{tab:ablation}) quantifies this: removing the reservoir /
water-value group costs $0.041$ micro AUC, while no other group costs
more than $0.004$. The final configuration adds the winning feature arm
from Table~\ref{tab:arms}.

\textbf{Hyperparameters.} Table~\ref{tab:xgbparams} lists the final
configuration. Values were selected by a closed screen scored on the 2022
validation year under the adoption bar: a candidate had to beat the
incumbent by $+0.0015$ micro AUC (the measured seed-to-seed noise floor
is $\pm 0.0014$) to be kept. Early stopping selects the number of trees
using only the last 90 days of the training period, never the scoring
year.

\begin{table}[h]
    \centering
    \begin{tabular}{lrl}
        \hline
        Parameter & Value & Search range \\
        \hline
        Learning rate & 0.05 & fixed \\
        Max depth & 7 & $\{6, 7, 8\}$ \\
        Min child weight & 5 & $\{5, 10\}$ \\
        Reg lambda & 10.0 & $\{1, 10, 30\}$ \\
        Reg alpha & 0.0 & fixed \\
        Subsample / colsample by tree & 0.8 / 0.8 & fixed \\
        Tree method & hist & fixed \\
        Max trees & 400 & early stopping, 25 rounds \\
        Eval metric & AUC & --- \\
        Seed & 0 & --- \\
        \hline
    \end{tabular}
    \caption{Final XGBoost configuration, shared by all 14 per-generator
    models. The search column gives the values tried in the closed
    hyperparameter screen.}
    \label{tab:xgbparams}
\end{table}

\textbf{Ensembling.} All ensembling is done on raw probabilities: the
rank-normalized variant collapses the metric (0.9556 vs.\ 0.9809,
Section~\ref{sec:metric}) and was rejected. A cross-family GBDT+GBDT
ensemble did not clear the adoption bar.
% TODO: if the final submission averages feature-set variants, name them
% here; the repo contains single-configuration submissions only.

\textbf{Results.} Table~\ref{tab:xgbprogress} summarizes the development
on the 2022 validation year. Leave-one-year-out cross-validation gives
pooled out-of-fold micro AUC 0.9758 over all 2922 cases, weakest fold
2020 at 0.9018. The selected Kaggle submission scores 0.99066 public.

\begin{table}[h]
    \centering
    \begin{tabular}{lrr}
        \hline
        Configuration & Val 2022 & Status \\
        \hline
        v2 baseline (193 features) & 0.9823 & reference \\
        + 12h LP allocation & 0.9842 & adopted \\
        + min-volume margin (final) & 0.9857 & adopted, submitted \\
        GNN + TCN (own attempt) & 0.9794 & rejected \\
        \hline
    \end{tabular}
    \caption{Development progression on the 2022 validation year. The
    sequence-model attempt is included for completeness; it scored below
    both tree anchors and was not adopted.}
    \label{tab:xgbprogress}
\end{table}

\textbf{What did not work.} (i) Rank-normalized ensembling ($-0.025$).
(ii) Cross-family GBDT+GBDT ensembling (below the adoption bar).
(iii) Five feature-search rounds --- inflow anomalies, extended price
history, and duplicate encodings of the same physics --- none of which
reached $+0.0015$ on 2022.

\subsection{Temporal Convolutional Network}
\label{sec:tcn}

\textbf{Setup.} One jointly trained network predicts all 14 generators at
all 168 hours: 343{,}034 parameters, consuming 913 hourly and 182 static
input channels. There is no ensemble and no pretrained model.

\textbf{Architecture.} A width-64 temporal path sees the complete
planning week: seven residual blocks, each with two kernel-3
convolutions, dilations $1,2,4,8,16,32,64$, giving a receptive field of
509 hours. Padding is symmetric (noncausal) because the complete
exogenous week is supplied at prediction time. Nine width-16 plant
branches receive explicitly routed local reservoir, tunnel and river
features and return per-generator corrections; their final layers start
at zero, so the branches can only learn to help. Everything trains
jointly under binary cross-entropy.

\textbf{Training.} AdamW (learning rate 0.001, weight decay 0.0001),
batch size 32 cases, dropout 0.1, gradient clipping 1.0, seed 42. Inputs
are asinh-encoded and scaled with statistics computed on fitting cases
only (clip at 10). Training budgets are fixed in advance: 20/42 epochs
for the two validation folds, 22 epochs for the final fit on all 2922
cases. A full fit takes about 500 seconds and 9.1 GiB on four CPU
threads.

\textbf{Results.} On the purged 2020 and 2021 folds the network reaches
micro AUC 0.9809 and 0.9893. The selected Kaggle submission scores
0.99261 public. Development was feature-driven: the terminal
water-value conversion moved the reported public score from 0.99197 to
0.99240, and refining the LP allocation interval from 12 to 6 hours
reached 0.99261. A matched comparison showed 6-hour intervals beat both
12-hour and 1-hour on both folds (1-hour cases also required artificial
storage slack in the LP).

\subsection{Comparison on the Same Folds}
\label{sec:comparison}

On the purged 2020/2021 folds, the sequence model scores
0.9809/0.9893. An earlier XGBoost generation with topology and
allocation features (298 features) scored 0.9830/0.9879 on the same
years: the tree model is better on 2020, the sequence model better on
2021. Neither family dominates; the final submissions therefore carry
one model of each family, and the private leaderboard decides.
% TODO: the fold construction of the two rows above is documented in
% docs/experiments.md; state any purging differences explicitly if the
% grader may ask. A naive-baseline row would strengthen this table.

\subsection{Unsuccessful Model Families}

\begin{itemize}
    \item \textbf{One pooled model for all generators} (trees): scored
    worse than the per-generator split.
    % TODO: insert the pooled score if recoverable, else keep qualitative.
    \item \textbf{GNN + TCN} (own sequence attempt with graph message
    passing): 0.9794 on val 2022, below both tree anchors (0.9823,
    0.9842); rejected.
    \item \textbf{Fourteen independent per-unit TCNs}: 0.9745/0.9816 on
    the 2020/2021 folds, below the shared network (0.9749/0.9841 in the
    same round); rejected.
    \item \textbf{Neural graph message passing in the shared TCN}: won
    both local years but scored lower publicly than local plant branches
    (0.99159 vs.\ 0.99176); rejected.
    \item \textbf{Cross-family GBDT+GBDT ensembling}: below the adoption
    bar.
    \item \textbf{Historical tree baselines}: CatBoost 0.9678/0.9827 and
    LightGBM 0.9644/0.9834 on the 2020/2021 folds; superseded.
\end{itemize}

\section{Model Interpretation}
\label{sec:interpretation}

\textbf{Which features matter.} The feature-group ablation
(Table~\ref{tab:ablation}) and the univariate screen
(Section~\ref{sec:individual}) agree: the reservoir / water-value group
dominates, and no small subset suffices. XGBoost's total-gain importance
rankings per generator confirm the same ordering.
% TODO: add the per-generator feature-importance figure (one bar panel),
% computed from the fitted trees -- no new training needed.

\textbf{Per-generator behavior.} Validation AUCs range from about 0.92
(Lio, Haukeli) to above 0.99 (Byrte, Vesle Kjela). The hard units fail
for different reasons: Haukeli sits at the end of a decade-long ON-rate
drift (0.87 in 2015 to 0.14 in 2022), while Lio responds to upstream
water context in the opposite direction of other units
(Section~\ref{sec:pairs}).
% TODO: add the per-unit AUC table/figure for the final configuration.

\textbf{Example weeks.} Predicted ON-probabilities track the
price-minus-water-value signal hour by hour, with the expected
stickiness: probabilities stay high through short unfavorable hours
rather than dropping, consistent with start/stop economics.
% TODO: add one example-week figure (price curve + predicted
% probabilities for 2-3 generators). Requires saved predictions.

\section{Limitations}

The labels come from an optimizer, not measurements, and the same
physical hour carries different labels in overlapping cases, capping the
achievable AUC (in-sample ceiling 0.9985). The LP water-system features
are a simplified relaxation, not the original scheduler: they do not
enforce binary commitment, changing-head optimization or spill behavior
exactly. The development years were reused across many comparisons and
are not pristine confirmation sets; 2022 in particular was scored
repeatedly during development. Public leaderboard scores are reported for
reproduction purposes only; grading uses the private leaderboard, whose
ranking is unknown to us. No external data was used.

\section{Use of AI}
% TODO: describe how AI tools were used in the project.

\bibliographystyle{alpha}
\bibliography{sample}

\end{document}